\section{Results}

\subsection{Main Results}\label{main-results}

Each task-topology pair ran three times with different random seeds, producing 738 primary runs (82 tasks $\times$ 3 topologies $\times$ 3 seeds). Table~\ref{tab:main_results} reports mean accuracy with 95\% confidence intervals.

\begin{table}[t]
\centering
\caption{Replicated results across 82 tasks. Accuracy (\%) is strict full-credit, averaged over 3 seeds. 95\% CIs from seed-level variance. All runs use Claude Opus 4.6.}
\label{tab:main_results}
\small
\begin{tabular}{lccc}
\toprule
Topology & Mean Accuracy & 95\% CI & Run 1 / 2 / 3 \\
\midrule
Flat & 28.9 & $\pm$2.1 & 29.3 / 26.8 / 30.5 \\
Hierarchical & 55.7 & $\pm$4.2 & 59.8 / 52.4 / 54.9 \\
Debate & 65.4 & $\pm$0.8 & 64.6 / 65.8 / 65.8 \\
\bottomrule
\end{tabular}
\end{table}

\begin{figure}[t]
\centering
\includegraphics[width=\columnwidth]{fig2_main_results.pdf}
\caption{Mean accuracy by topology across 82 tasks (3-seed average). Error bars show 95\% confidence intervals. Debate leads overall, but the advantage is task-dependent (see Figure~\ref{fig:crossover}).}
\label{fig:main_results}
\end{figure}

CIs are non-overlapping between flat and both multi-agent topologies. The hierarchical-debate gap has wider CIs due to run-to-run variance in hierarchical performance (52.4--59.8\%), but the ordering never reverses in any single run. For practitioners, this means the question is not whether to use multi-agent orchestration on hard tasks, but which multi-agent topology to use. A secondary backbone check on Claude Sonnet 4.6 (12 easy tasks, 36 runs) confirms the ceiling effect: all three topologies score 83.3\% on easy tasks, confirming that topology differentiation requires hard tasks.

\subsection{Statistical Rigor}\label{statistical-rigor}

We report effect sizes and corrected significance tests to distinguish real effects from noise.

\begin{table}[t]
\centering
\caption{Pairwise comparisons with effect sizes and corrected $p$-values. Cohen's $d$ computed from seed-level means. $P$-values from paired permutation tests with Holm-Bonferroni correction for 3 comparisons.}
\label{tab:effect_sizes}
\small
\begin{tabular}{lcccc}
\toprule
Comparison & $\Delta$ (pp) & Cohen's $d$ & Holm-Bonf.\ $p$ & Practical \\
\midrule
Debate vs.\ Flat & +36.5 & 8.72 & $<$0.001 & Large \\
Hier.\ vs.\ Flat & +26.8 & 4.81 & $<$0.001 & Large \\
Debate vs.\ Hier. & +9.7 & 2.14 & 0.016 & Med--Large \\
\bottomrule
\end{tabular}
\end{table}

All three pairwise comparisons survive Holm-Bonferroni correction at $\alpha = 0.05$. The effect sizes are large by conventional standards ($d > 0.8$). The debate-vs-hierarchical comparison has the smallest effect ($d = 2.14$), reflecting the higher variance in hierarchical performance across seeds. This variance is itself informative: hierarchical topology is more sensitive to random seed than debate, likely because the planner's initial decomposition commits the entire execution to a single strategy.

Bootstrap confidence intervals (10,000 resamples of task-level results) for the DIT routing classifier: 90\% accuracy, 95\% CI [84\%, 95\%]. The lower bound of 84\% still exceeds the best single-topology strategy (debate at 65.4\%) by nearly 20 percentage points.

Statistical significance does not automatically imply practical significance. For the flat-vs-multi-agent comparisons, the practical impact is unambiguous: a 27--37pp accuracy gap changes deployment outcomes. For the debate-vs-hierarchical comparison, the 9.7pp aggregate gap understates the practical impact, because the advantage is concentrated on specific task profiles (see Section~\ref{difficulty-dependent-topology-advantage}). On those profiles, the gaps reach 33--63pp.

\subsection{Difficulty-Dependent Topology Advantage}\label{difficulty-dependent-topology-advantage}

Of the 70 hard tasks, hierarchical was the sole winner on 24 and debate on 29. The split tracks DIT dimensions. Hierarchical's wins cluster on high-D tasks ($D \geq 0.6$, $n{=}24$, mean $D{=}0.82$): when a task decomposes into independent subtasks, the planner assigns one per executor. Averaged across three seeds, hierarchical reaches 75\% on these tasks (Figure~\ref{fig:crossover}). Debate's wins cluster on high-I tasks ($I \geq 0.6$, $n{=}22$, mean $I{=}0.85$): two agents produce complementary partial solutions. On high-I tasks, debate reaches 77\%.

\begin{figure}[t]
\centering
\includegraphics[width=\columnwidth]{fig3_crossover.pdf}
\caption{Difficulty-dependent crossover: hierarchical dominates on high-D tasks while debate dominates on high-I tasks. Each bar shows 3-seed mean accuracy with 95\% CI.}
\label{fig:crossover}
\end{figure}

The crossover is sharp. On high-D tasks, hierarchical outperforms debate by 33pp (75\% vs.\ 42\%). On high-I tasks, debate outperforms hierarchical by 63pp (77\% vs.\ 14\%). Each topology fails in the other's territory for predictable structural reasons: hierarchical fragments cross-cutting concerns across executors, while debate's independent solvers produce divergent designs that resist clean merging.

\subsection{DIT Dimension Importance}\label{dit-dimension-importance}

Which dimension matters most for topology selection? The answer depends on the task category.

\begin{table}[t]
\centering
\caption{Predictive power of each DIT dimension for topology selection, measured by point-biserial correlation between dimension value and correct-topology prediction. Higher $|r|$ indicates the dimension is more diagnostic.}
\label{tab:dit_importance}
\small
\begin{tabular}{lcccc}
\toprule
Category & $r$(D) & $r$(I) & $r$(T) & Dominant \\
\midrule
Coding & 0.81 & $-$0.62 & 0.23 & D \\
Reasoning & $-$0.41 & 0.87 & 0.18 & I \\
Research & 0.54 & $-$0.38 & 0.47 & D (T sec.) \\
\midrule
\textbf{Overall} & \textbf{0.68} & $\mathbf{-0.59}$ & \textbf{0.31} & \textbf{D} \\
\bottomrule
\end{tabular}
\end{table}

Decomposability is the strongest single predictor overall ($r = 0.68$), but iterativeness dominates for reasoning tasks ($r = 0.87$). Tool diversity (T) is a secondary signal that becomes relevant for research tasks requiring cross-tool coordination. For practitioners, a simple rule follows: check decomposability first. If $D \geq 0.6$, use hierarchical. If D is low and $I \geq 0.6$, use debate. The full DIT distance metric refines this heuristic for tasks in the ambiguous middle region ($0.3 < D, I < 0.6$), where 6 of the 8 routing mismatches occur.

\subsection{Agent-Count Control and Token Efficiency}\label{agent-count-control}

Multi-agent topologies use more compute than flat. To isolate structure from compute, we gave a single flat agent 2$\times$ the token budget on 20 hard tasks. Flat-2$\times$ reaches 35\%, up from 20\% at standard budget. Multi-agent topologies score 65\% on the same tasks at comparable cost. The decomposition: of the 45pp gap between flat and multi-agent, roughly 15 points come from additional compute and 30 from topology structure. \textbf{Structure accounts for two-thirds of the multi-agent advantage.}

\begin{figure}[t]
\centering
\includegraphics[width=\columnwidth]{fig5_agent_control.pdf}
\caption{Agent-count control experiment on 20 hard tasks. Flat-2$\times$ (doubled token budget) partially closes the gap with multi-agent topologies, but structure accounts for ${\sim}$2/3 of the advantage.}
\label{fig:agent_control}
\end{figure}

Simply giving a single agent more tokens recovers only one-third of the multi-agent benefit. The remaining two-thirds requires the organizational structure itself --- role differentiation, decomposition planning, or adversarial critique --- functions that emerge from topology, not compute.

On hard tasks, flat is cheapest in raw tokens (${\sim}$2,490 vs.\ ${\sim}$4,770 hierarchical, ${\sim}$4,710 debate) but most expensive per correct answer (${\sim}$13,400 tokens/success vs.\ ${\sim}$9,000 hierarchical and ${\sim}$8,000 debate). For budget-constrained deployments, the token-per-success metric inverts the naive cost ranking: debate is 40\% cheaper than flat per unit of correct output.

\subsection{DIT Routing Accuracy}\label{dit-routing-accuracy}

The routing classifier selects the best topology with 90\% accuracy (74/82 tasks), against a 33\% random baseline. The 8 mismatches cluster where D and I scores are both moderate (0.3--0.6), the ambiguous region where no topology has a strong structural advantage. On tasks with extreme profiles ($D \geq 0.8$ or $I \geq 0.8$), routing accuracy reaches 100\% (36/36 tasks).

\subsection{Failure Mode Taxonomy}\label{failure-mode-taxonomy}

Why does each topology fail when mismatched? We analyzed all 34 topology-specific failures across the 70 hard tasks and identified three dominant failure modes per topology (full analysis in Appendix~\ref{appendix:failure-analysis}).

\begin{table}[t]
\centering
\caption{Failure mode taxonomy. Each topology has a characteristic failure pattern mapping to specific DIT regions.}
\label{tab:failure_modes}
\small
\begin{tabular}{p{1.4cm}p{2.0cm}cp{1.3cm}}
\toprule
Topology & Failure Mode & Freq & DIT Region \\
\midrule
Flat & Context overload & 7/18 & All hard \\
Flat & Anchoring bias & 4/18 & All hard \\
Flat & Pipeline shortcut & 4/18 & All hard \\
Hier. & Bad decomposition & 9/9 & High-I \\
Debate & Merge inconsist. & 7/7 & High-D \\
\bottomrule
\end{tabular}
\end{table}

The symmetry is striking: hierarchical's failure mode (bad decomposition on high-I tasks) is the structural mirror of debate's failure mode (merge inconsistency on high-D tasks). Each topology's weakness is the other's strength. This follows directly from their positions in DIT space. Hierarchical assumes tasks decompose cleanly (high-D prior), so it fails when they do not. Debate assumes tasks benefit from adversarial refinement (high-I prior), so it fails when tasks need coherent modular construction instead.

Flat's failure modes, by contrast, are not DIT-specific. Context overload, anchoring, and shortcutting are general single-agent limitations that affect all hard tasks. This explains why flat's accuracy on hard tasks (18.6\%) is uniformly low regardless of DIT profile.

\subsection{Limitations}\label{sec:results-limitations}

We want to be direct about what 82 tasks can and cannot establish. 794 total runs with three-seed replication provide CIs for primary comparisons, but 82 tasks may not generalize to all domains. Three topologies do not exhaust the design space --- the characterized-but-unevaluated role-playing and RL-orchestrated topologies could occupy different failure-mode niches. DIT scores were author-assigned ($\kappa \geq 0.79$, scored before experiments); third-party annotation is needed. The primary backbone is Claude Opus 4.6; topology-model interactions across families remain untested. Hand-constructed tasks introduce potential authorship bias; all are released for verification. The error bars exist, the ordering is stable across seeds, and the direction is clear: topology choice is invisible on easy tasks, decisive on hard ones, and predictable from DIT characteristics.
